\scanpage{36}
% archive sheet number: 35
On a $A\otimes_{\mathbb Z_p}W$ \unsure{étale} sur $A$, ext. \unsure{résiduelle|résiduels}
$\begin{aligned}[t] k\otimes_{\mathbb F_p}k&\simeq k^\Sigma,\quad\Sigma=\operatorname{Aut} k\\[-2pt]&\hphantom{\simeq k^\Sigma,\quad\Sigma}=\operatorname{Aut} W\end{aligned}$
\[
A\otimes_{\mathbb Z_p}W\simeq A^\Sigma,\qquad \Sigma\simeq\mathbb Z/f\mathbb Z\quad\text{en correspondance}
\]
\hfill bijection canonique avec

\hfill $[0,f-1]$ par $\unsure{\nu|v}\longmapsto\sigma^\nu$

\hfill \deleted{$(\unsure{\sigma^\nu=\mathrm F^p|\sigma=\mathrm F^p})$}

On a
\[
\sigma_A(x_0,x_1,\ldots,x_{f-1})=(x_1,\ldots,x_{f-1},x_0)
\]

La donnée d'un module \inserted{libre} $M$ \inserted{sur $A\otimes_{\mathbb Z_p}W$} équivaut à la donnée de modules $(\deleted{\unsure{A}}M_\nu)_{\nu\in\Sigma}$ sur $A$ ; le rang de $\deleted{\unsure{A}}M$ est la somme des rangs des $M_\nu$. \inserted{(on le veut égal à $h$)} La donnée d'un \unsure{endom} $F_M$ $\sigma_A$-linéaire équivaut à la donnée d'hom. $A$-linéaires
\diagram{\begin{tikzcd}[ampersand replacement=\&,column sep=small]
M_0\arrow[r,"\varphi_0"] \& M_1\arrow[r,"\varphi_1"] \& \cdots\arrow[r] \& M_{f-1}\arrow[r,dash,"\varphi_{f-1}"] \& M_0
\end{tikzcd}}

La condition $(*)$ \deleted{\unsure{s'écrit}} implique que les $f$ composés circulaires sont des isogénies, donc les $\varphi_i$ tous injectifs, donc les $\varphi_i$ des isogénies (car les rangs des $M_\nu$ tous égaux \ldots). \deleted{Donc} --- Utilisant maintenant le fait que la somme des rangs des $M_\nu$ est $h$, on trouve que les $M_\nu$ sont inversibles.
% slanted
\margin{ce serait aussi une \deleted{conséquence de $(*1)$} \deleted{\unsure{cas}} de $(*)$ avec $\pi$ un élément non nul \uline{quelconque} de $A$}

\deleted{On peut alors trouver un unique module inversible}

On peut alors identifier tous les $M_\nu$ à \deleted{$M_0=\Omega$,} \deleted{$\Omega$ sur $A$,} $M_0=t$, \deleted{\unsure{qu'on a|qui a i}} en choisissant \deleted{\unsure{comme|can.}} pour \unsure{tt|$M$} $\unsure{\nu|\nu'}$,
\inserted{$0\leq\nu\leq f-1$,} l'isom $M_0\xrightarrow{\sim}M\unsure{.|\!{}_\nu|\!{}_{i}}$ donné par
$\deleted{\unsure{p}}\pi^{-m_\nu}(\varphi_{\nu-1}\varphi_{\nu-2}\ldots\varphi_0)$, où $m_\nu$ est un entier $\geq 0$ uniquement déterminé ($\unsure{0\leq|0=}m_0\leq m_1\leq\ldots\leq m_{f-1}$). On trouve donc un isom can.
\[
 M\simeq\Omega\otimes_{\mathbb Z_p} W=\Omega\otimes_A(A\otimes_{\mathbb Z_p}W)
\]
